\documentclass[11pt]{article}
\usepackage{fontspec}
\setmainfont{Times New Roman}
\setsansfont{Arial}
\setmonofont{Consolas}
\usepackage{amsmath,amssymb}
\usepackage{geometry}
\usepackage{microtype}
\geometry{a4paper,margin=3cm}
\setlength{\parindent}{1.25em}
\setlength{\parskip}{0pt}
\emergencystretch=30em
\allowdisplaybreaks
\raggedbottom
\newcommand{\ideal}[1]{\mathfrak{#1}}
\newcommand{\eqzero}[1]{\equiv 0\,(#1)}

\begin{document}

% Унит: Папер 19, Сецтион 10. Дирецт Интерславиц Латин транслатион фром те Герман финал-аудитед соурце слице.

\setcounter{footnote}{38}
\subsection*{\S~10. Специалны случај полиномној области.}

1. Нехај основна колцева област \(\Sigma\) состоји из всех полиномов од \(x_1,
\ldots,x_n\) с либоволными комплексными коефициентами; за ньу, по Хилбертовом теорему о модулном базису (Ann. 36), изполньено јест условје конечности. Иде о связ наших теоремов с познаными теоремами теорије елиминације и теорије модулов.

Тутој связ устанавльа следечи специалны случај познаного Хилбертового теорема:\footnote{О полных системах инвариантов. Math. Ann. 42 (1893), \S~3, с. 313.}

Ако \(f\) изчезаје за всако конечно система вредностиј \(x_1,
\ldots,x_n\), кторо јест обча нула всех полиномов простого идеала \(\ideal P\) (нула од \(\ideal P\)), тогда \(f\) јест дельива чрез \(\ideal P\). Другими словами: просты идеал \(\ideal P\) состоји из целости полиномов, кторе изчезајут в јего нулах.\footnote{Како показал Lasker [Math. Ann. 60 (1905), с. 607], тутој специалны случај можно в случају хомогеных форм доказати и директно; обратно, из того знову следује Хилбертов теорем в хомогеном и нехомогеном случају, кторы можно изректи тако: ако идеал \(\ideal R\) изчезаје во всех нулах од \(M\), тогда нека потенца од \(\ideal R\) јест дельива чрез \(M\). Але тој теорем, и такоже специалны случај, важи само ако вредностны обсег од \(x\) јест алгебраично замкньены; зато не може следовати само из наших теоремов, але мора употребити ексистенцију коренов, напр. на том месту, кде идеал, чије нуле состоје само из \(x_1=0,
\ldots,x_n=0\), всебује все степенове продукты од \(x\) од неког степена надалье. Остаток доказа може се, с употребеньем наших теоремов, неколико упростит взходно к Lasker. Lasker бо мора употребити Хилбертов теорем такоже за доказ разложеньа идеала на највечше примарне.}

Ако зато продукт \(fg\) изчезаје за все нуле од \(\ideal P\), тогда изчезаје најмене једен фактор; нуле творут неразложиму алгебраичну многовидност. Ако се обратно та дефиниција неразложиме многовидности положи в основу, показује се, же целост полиномов изчезајучих в неразложимој многовидности твори просты идеал; просты идеал и неразложима многовидност зато одповедајут једнозначно једен другому. Због \(\ideal P^\rho\eqzero{\ideal Q}\), \(\ideal Q\eqzero{\ideal P}\) нуле примарного идеала совпадајут с нулами јего придруженого простого идеала.\footnote{То својство примарного идеала, именно имети неразложиму многовидност, Macaulay (срав. Увод) бере за дефиницију, меджутым как Lasker в дефиницију вкльучаје само појем многовидности твора, а в осталом дефинира абстрактно. Примарне идеалы изчезајуче само за \(x_1=0,
\ldots,x_n=0\) имајут у Lasker особно положенье.} Представјенье идеала како најменшего обчего многократного највечших примарных идеалов даје зато разложенье всех нул идеала на неразложиме многовидности; и, како показал Lasker, важи такоже обратно. Доказ једнозначности придруженых простых идеалов одповедаје тут фундаменталному теорему теорије елиминације о једнозначној разложимости алгебраичне многовидности на неразложиме многовидности; за специалне полиномне области, кде не обстаја једнозначно продуктно представјенье полиномов чрез неразложиме полиномы области и следовательно не обстаја ни теорија елиминације, може он важити како еквивалент тога теорема теорије елиминације.

Неразложиме многовидности одповедајуче изолованым примарным идеалам сут точно оне, кторе выступајут в ``минималној резолвенте'';\footnote{Срав. напр. J. König, Einleitung in die allgemeine Theorie der algebraischen Größen (Leipzig, Teubner, 1903), с. 235.} бо нуле всакого неизолованого највечшего примарного идеала сут Једночасно нуле најмене једного изолованого, именно такого, чиј придружены просты идеал јест дельив чрез просты идеал неизолованого идеала. Зато једнозначност изолованых примарных идеалов даје в експонентах нове инвариантне числа многократности. Такоже теоремы једнозначности при разложеньу примарных идеалов на неразложиме можно сматрјати како дополньенье теорије елиминације в смыслу многократности.

По тых заметах различна разложеньа можно толковати по јих односу к алгебраичным многовидностјам. Попарно взаимно простым идеалам одповедајут таке многовидности, кторе не имајут обчеј нулы; при взаимно релативно простых идеалах ниједна неразложима многовидност једного идеала не јест Једночасно обча нула; највечше примарне идеалы изчезајут само на неразложимых многовидностјах, все различных меджу собоју; при разложеньу на неразложиме идеалы те саме неразложиме многовидности могут выступати такоже многократно.

Нехај буде ешче замечено, же наместо обчеј полиномној области можно положити в основу и област всех хомогеных форм, бо легко се уверити, же и при там важечих связах -- сложенье јест дефинирано само за формы једнакого степена -- обче теоремы оставајут в силе.\footnote{Же тут пак в случају многозначности поред хомогеных могут ексистовати и нехомогене разложеньа, показује примјер \( (x^3,xy,y^3)=[(x^3,y),(y^3,x)]=[(xy,x^3,y^3,x+y^3),(xy,x^3,y^3,y+x^2)]\).}

Најпростши примјер четырих различных разложень -- формулы за него следујут ниже -- по горньем јест дан, напр., прјамоју и двема к ньеј мимобежныма прјамыма, кторе се пересекајут меджу собоју, из кторых једна всебује точку, различну од точкы пересека, во вышеј многократности. Разложеньу на взаимно просте неразложиме идеалы одповедаје разложенье на прјаму и к ньеј мимобежну многовидност; та многовидност при разложеньу на релативно просте неразложиме идеалы разпада се на две прјамы; разложеньу на највечше примарне идеалы одповедаје оддельенье точкы вышеј многократности, меджутым как разложенье на неразложиме идеалы условльује разложенье тој точкы.

Ако се та точка взме за почеток, прјама проходеча чрез ньу за ос \(y\), прјама пересекајуча ју паралелно к оси \(x\), а к ньеј мимобежна прјама паралелно к оси \(z\), тогда така конфигурација јест, напр., представльена следечими неразложимыми идеалами:\footnote{Из ньих прве три сут неразложиме како просте идеалы; \(\ideal B_4\), понеже има само делителье \((x^2,y,z)\) и \((x,y,z)\); \(\ideal B_5\), понеже всакы делитель всебује полином \(xy\).}
\[
\begin{gathered}
 \ideal B_1=(x-1,y),\quad
 \ideal B_2=(y-1,z),\quad
 \ideal B_3=(x,z),\quad
 \ideal B_4=(x^3,y,z),\quad
 \ideal B_5=(x^2,y^2,z).
\end{gathered}
\]
Придружене просте идеалы бывајут:
\[
 \ideal P_1=\ideal B_1,
 \quad \ideal P_2=\ideal B_2,
 \quad \ideal P_3=\ideal B_3,
 \quad \ideal P_4=\ideal P_5=(x,y,z).
\]
Највечше примарне идеалы бывајут:
\[
 \ideal Q_1=\ideal B_1,
 \quad \ideal Q_2=\ideal B_2,
 \quad \ideal Q_3=\ideal B_3,
 \quad \ideal Q_4=[\ideal B_4,\ideal B_5]=(x^3,y^2,x^2y,z).
\]
Релативно просте неразложиме идеалы бывајут:
\[
 \ideal R_1=\ideal Q_1,
 \quad \ideal R_2=\ideal Q_2,
 \quad \ideal R_3=[\ideal Q_3,\ideal Q_4]=(x^3,x^2y,xy^2,z).
\]
Взаимно просте неразложиме идеалы бывајут:
\[
 \ideal S_1=\ideal R_1,
 \quad
 \ideal S_2=[\ideal R_2,\ideal R_3]=((y-1)x^3,(y-1)x^2y,(y-1)xy^2,z).
\]
Из того приходи целы идеал:
\[
\begin{aligned}
 \ideal M&=[\ideal S_1,\ideal S_2]=\ideal S_1\ideal S_2 \\
 &=((x-1)(y-1)x^3,(y-1)x^2y,(y-1)xy^2,(x-1)z,y(y-1)x^3,yz),
\end{aligned}
\]
при чем
\[
 1=-(y-1)x^3+(y-1)(x^3-1)+y,
\]
\[
 (y-1)(x^3-1)+y\eqzero{\ideal S_1},
 \qquad -(y-1)x^3\eqzero{\ideal S_2}.
\]
Притом идеалы \(\ideal B_1\), \(\ideal B_2\), \(\ideal B_3\) сут изоловане, зато такоже при разложеньах на неразложиме и на највечше примарне идеалы једнозначно определьене. Идеалы \(\ideal B_4\) и \(\ideal B_5\), односно \(\ideal Q_4\), сут неизоловане; оне не сут једнозначно определьене, але напр. замениме чрез
\[
 \ideal D_4=(x^3,y+\lambda x^2,z),
 \qquad
 \ideal D_5=(x^2+\mu xy,y^2,z),
\]
такоже \(\ideal Q_4\) јест заменимы чрез
\[
 \overline{\ideal Q}_4=(x^3,x^2y,y^2+\lambda xy,z).
\]

2. Тако само како обча (и цело-числова) полиномна област, такоже всака конечна област целости из полиномов -- како показује Хилбертов теорем о модулном базису -- изполньа условје конечности,\footnote{Обратно, ако за полиномну област изполньено јест условје конечности и всакы полином допушчаје најмене једно представјенье, кде множительи сут нижшего степена по \(x\), тогда област јест конечна област целости.} при чем коефициенты могут быти придане либоволному польу. Дајемо ешче примјер за област всех парных полиномов, како најпростшу област, кде због \(x^2y^2=(xy)^2\) не обстаја једнозначно продуктно представјенье полиномов чрез неразложиме полиномы области. Иде о ту саму конфигурацију како в горньем примјеру, ктора тепер јест дана следечими неразложимыми идеалами:
\[
\begin{gathered}
 \ideal B_1=(x^2-1,xy,y^2,yz),\quad
 \ideal B_2=(y^2-1,xz,yz,z^2),\\
 \ideal B_3=(x^2,xy,xz,yz,z^2),\quad
 \ideal B_4=(x^4,x^3y,y^2,xz,yz,z^2),\\
 \ideal B_5=(x^2,y^2,xz,yz,z^2).
\end{gathered}
\]
Придружене просте идеалы бывајут:
\[
 \ideal P_1=\ideal B_1,
 \quad \ideal P_2=\ideal B_2,
 \quad \ideal P_3=\ideal B_3,
 \quad \ideal P_4=\ideal P_5=(x^2,xy,y^2,xz,yz,z^2).
\]
Же \(\ideal P_1\) бываје просты идеал, следује из того, же всакы полином области достава форму
\[
 f=\varphi(z^2)+xz\,\psi(z^2)\pmod{\ideal P_1}.
\]
Нехај зато
\[
 f_1=\varphi_1(z^2)+xz\,\psi_1(z^2),
 \qquad
 f_2=\varphi_2(z^2)+xz\,\psi_2(z^2),
\]
тогда из \(f_1f_2\eqzero{\ideal P_1}\) следује важенье једначин
\[
 \varphi_1\varphi_2+z^2\psi_1\psi_2=0,
 \qquad
 \varphi_2\psi_1+\varphi_1\psi_2=0,
\]
и из того \(f_1\eqzero{\ideal P_1}\) или \(f_2\eqzero{\ideal P_1}\). Точно тако показује се, же \(\ideal P_2\) бываје просты идеал; \(\ideal P_3\) бываје просты идеал, бо всакы полином области модуло \(\ideal P_3\) бываје конгруентны полиному од \(y^2\); \(\ideal P_4\) состоји из всех полиномов области. Зато такоже \(\ideal B_1\), \(\ideal B_2\) и \(\ideal B_3\) сут неразложиме како просте идеалы; \(\ideal B_4\) и \(\ideal B_5\) пак имајут всака само једины властны делитель \(\ideal P_4\), сут зато нужно неразложиме.

Из неразложимых идеалов произходет највечше примарне:
\[
\begin{gathered}
 \ideal Q_1=\ideal B_1,
 \quad \ideal Q_2=\ideal B_2,
 \quad \ideal Q_3=\ideal B_3,
 \quad
 \ideal Q_4=[\ideal B_4,\ideal B_5]=(x^4,x^3y,y^2,xz,yz,z^2),
\end{gathered}
\]
релативно просте неразложиме идеалы:
\[
 \ideal R_1=\ideal Q_1,
 \quad \ideal R_2=\ideal Q_2,
 \quad
 \ideal R_3=[\ideal Q_3,\ideal Q_4]=(x^4,x^3y,x^2y^2,xy^3,xz,yz,z^2),
\]
взаимно просте неразложиме идеалы:
\[
 \ideal S_1=\ideal R_1,
 \quad \ideal S_2=[\ideal R_2,\ideal R_3]
 =((y^2-1)x^4,(y^2-1)x^3y,(y^2-1)x^2y^2,
 (y^2-1)xy^3,xz,yz,z^2).
\]
Како в првом примјеру, бываје:
\[
 [\ideal B_4,\ideal B_5]=[\ideal D_4,\ideal D_5],
 \qquad
 [\ideal Q_3,\ideal Q_4]=[\ideal Q_3,
\overline{\ideal Q}_4],
\]
кде
\[
 \ideal D_4=(x^4,xy+\lambda x^2,\ldots),\qquad
 \ideal D_5=(x^2+\mu xy,\ldots),\qquad \lambda\mu\ne1,
\]
и
\[
 \overline{\ideal Q}_4=(x^4,x^3y,y^2+\lambda xy,\ldots).
\]
Остале неразложиме односно највечше примарне идеалы \(\ideal B_1\), \(\ideal B_2\), \(\ideal B_3\) сут како изоловане идеалы једнозначно определьене.

\end{document}

